% TOP_PROBLEM: {"id": 30007321, "title": "Computation by a fixed particle with eventually stationary forcing", "category_id": 9, "category": {"id": 9, "name": "pde", "display_name": "Partial Differential Equations", "description": "PDEs and their applications in physics and geometry.", "slug": "pde", "order_index": 9, "created_at": "2026-07-31T15:26:25.671Z"}, "status": "open"}
% FIRST_POSTED: 2026-10-06
% ATTEMPT_STATUS: solved
% OPENAI_CLAIM: {"schema_version":1,"source_repo":"https://github.com/openai/math","source_commit":"adc7f1241b42e322a6451854ab7e4b4c146bf78a","release_date":"2026-10-06","match":"full","resolution":"proved","independently_reviewed":false,"second_pass":"agreed","families":[{"family":"376","title":"Universal computation in forced Navier–Stokes flows","subject":"Partial differential equations","manuscripts":[{"title":"Universal Computation with Eventually Stationary Navier–Stokes Forcing","pdf_path":"preprints/Universal-Computation-with-Eventually-Stationary-Navier-Stokes-Forcing-September-27-2026/manuscript.pdf","readme_path":"preprints/Universal-Computation-with-Eventually-Stationary-Navier-Stokes-Forcing-September-27-2026/README.md","date":"2026-09-27","theorem_ref":"preprints/Universal-Computation-with-Eventually-Stationary-Navier-Stokes-Forcing-September-27-2026/build/main.tex:146 (thm:main)"}],"lean":{"doc_path":"lean/docs/376.md","comparators":[],"declarations":[],"covers_main_theorem":false},"reasoning_trace":null}]}
\documentclass[11pt]{article}
\usepackage[margin=25mm]{geometry}
\usepackage{fontspec}
\setmainfont{Latin Modern Roman}
\usepackage{amsmath,amssymb,mathtools}
\usepackage{unicode-math}
\setmathfont{Latin Modern Math}
\usepackage{enumitem}
\usepackage{xurl,hyperref,bookmark}
\hypersetup{colorlinks=true,urlcolor=blue}
\setcounter{secnumdepth}{0}
\setlength{\parindent}{0pt}
\setlength{\parskip}{5pt}
\setlength{\emergencystretch}{3em}
\begin{document}
\section*{Computation by a fixed particle with eventually stationary forcing}

\subsection{OpenAI's claim}
\textit{Eventually stationary forced computation}: The manuscript claims smooth mean-zero forcing on the flat three-torus that becomes stationary after time one and lets a fixed particle detect arbitrary machine halting. The force has a finite effective description and the fluid begins at rest.

\subsection{Outline of the argument}
\textit{Eventually stationary forced computation}: A reversible recording table is encoded by gapped tape rectangles. Ordered extraction into a parking region, reciprocal deformation, and delivery implement every branch without first running the computation. A third spatial coordinate drives the finite rule table autonomously; a fixed startup ramp changes rest into the autonomous field, after which the force is stationary.

\subsection{Scope relative to this problem}
\textit{Eventually stationary forced computation}: The computational rule table supplies a finite effective force prescription, which is stationary after the loading interval. Bounded derivatives, classical uniqueness and uniformly bounded kinetic energy accompany the particle event; no trajectory-dependent choice of forcing is made.

With the hypotheses and conventions above, the claimed result addresses the whole mathematical assertion on this page.

\subsection{Formal verification}
Family 376 describes Lean coverage for the selected sheet-program, rapidly decaying particle-detector, velocity-detector, and full-box transport statements described separately in its Comparator files.

\textit{Eventually stationary forced computation}: No manuscript-specific Comparator is listed in the pinned inventory.

The documented selection does not cover the entire manuscript assertion relevant here.

This site has not run the Lean development. These links describe OpenAI’s selected formal statements and should be read with their stated scope.

\subsection{Status caveat}
OpenAI reports that an internal model produced the manuscripts. Their arguments have not been independently verified as proofs by this site. Any portion outside the documented formal scope may contain mistakes, and the catalogue label reports the announced mathematical claim.

\subsection{Sources}
\begin{itemize}
\item OpenAI, \emph{Universal Computation with Eventually Stationary Navier–Stokes Forcing}, 2026-09-27; citation key \texttt{OAI:Universal-Computation-with-Eventually-Stationary-Navier-Stokes-Forcing-September-27-2026}. \href{https://github.com/openai/math/blob/adc7f1241b42e322a6451854ab7e4b4c146bf78a/preprints/Universal-Computation-with-Eventually-Stationary-Navier-Stokes-Forcing-September-27-2026/manuscript.pdf}{PDF (GitHub)}; \href{https://github.com/openai/math/raw/adc7f1241b42e322a6451854ab7e4b4c146bf78a/preprints/Universal-Computation-with-Eventually-Stationary-Navier-Stokes-Forcing-September-27-2026/manuscript.pdf}{PDF (direct)}; \href{https://github.com/openai/math/blob/adc7f1241b42e322a6451854ab7e4b4c146bf78a/preprints/Universal-Computation-with-Eventually-Stationary-Navier-Stokes-Forcing-September-27-2026/README.md}{Manuscript page and citation}; \href{https://github.com/openai/math/blob/main/preprints/Universal-Computation-with-Eventually-Stationary-Navier-Stokes-Forcing-September-27-2026/manuscript.pdf}{Latest version}.
\item Family 376: \href{https://github.com/openai/math/blob/adc7f1241b42e322a6451854ab7e4b4c146bf78a/lean/docs/376.md}{Lean scope}.
\item \href{https://github.com/openai/math/blob/adc7f1241b42e322a6451854ab7e4b4c146bf78a/CONTENTS.md}{OpenAI catalogue entry}.
\end{itemize}
\end{document}
